\documentclass[conference]{IEEEtran}
\IEEEoverridecommandlockouts
\usepackage[T1]{fontenc}
\usepackage[utf8]{inputenc}
\usepackage[spanish,es-nolayout,es-tabla,es-nodecimaldot,es-noshorthands]{babel}
\usepackage[nocompress]{cite}
\usepackage{amsmath,amssymb}
\usepackage{graphicx}
\usepackage{textcomp}
\usepackage{xcolor}
\usepackage{array}
\usepackage{stfloats}
\usepackage{tikz}
\usetikzlibrary{positioning,arrows.meta,calc,fit}
\usepackage{url}
\def\UrlBreaks{\do\/\do-\do.\do_\do?\do=\do\&\do\%}

% Nombres de los componentes en español con la forma de IEEE.
\addto\captionsspanish{%
  \renewcommand{\figurename}{Fig.}%
  \renewcommand{\tablename}{TABLA}%
  \renewcommand{\abstractname}{Resumen}%
  \renewcommand{\refname}{Referencias}%
}
\renewcommand{\IEEEkeywordsname}{Palabras clave}
\newcommand{\code}[1]{\texttt{#1}}
\urldef{\urlbalogh}\url{https://llvm.org/devmtg/2018-04/slides/Balogh-Finding%20Iterator-related%20Errors%20with%20Clang%20Static%20Analyzer.pdf}

\begin{document}

\title{Detección de iteradores, referencias y punteros invalidados en contenedores de C++ con un checker para Clang Static Analyzer derivado de la norma}

\author{\IEEEauthorblockN{Jason Sebastián Rodríguez Chávez}
\IEEEauthorblockA{\textit{Escuela Profesional de Ingeniería de Sistemas} \\
\textit{Universidad Católica de Santa María}\\
Arequipa, Perú}
\and
\IEEEauthorblockN{Mayerly Katrina Aranibar Zegovia}
\IEEEauthorblockA{\textit{Escuela Profesional de Ingeniería de Sistemas} \\
\textit{Universidad Católica de Santa María}\\
Arequipa, Perú}
}

\maketitle

\begin{abstract}
Usar un iterador, una referencia o un puntero a un elemento de un contenedor de la biblioteca estándar de C++ después de una operación que lo invalida es comportamiento indefinido. Las reglas que deciden cuándo ocurre dependen del contenedor, de la categoría del identificador, de la posición del elemento y, en algunos casos, del estado del programa; los analizadores estáticos las aproximan con listas de operaciones o se limitan a los iteradores. Este artículo presenta un checker para Clang Static Analyzer que aplica las reglas de invalidación de C++17, escritas como datos y con su cláusula, sobre la ejecución simbólica del analizador. El modelo distingue tres veredictos: uso inválido confirmado, uso posiblemente inválido cuando la invalidación depende de la capacidad o del factor de carga, y uso válido; además, reconoce las garantías que el propio código establece. Se evaluó sobre un banco de 750 casos generado a partir de la misma norma y contrastado con dos instrumentos dinámicos en 852 ejecuciones sin contradicciones. La configuración de referencia del analizador detecta 80 de 342 usos inválidos y ninguno sobre referencias o punteros retenidos; con el checker se detectan los 342, sin falsas alarmas en 306 casos seguros y con un sobrecosto mediano de tiempo de 0.3 por ciento. En 781 unidades de traducción de cuatro proyectos de código abierto no produjo avisos confirmados falsos.
\end{abstract}

\begin{IEEEkeywords}
análisis estático, C++, detección de defectos, ejecución simbólica, invalidación de iteradores, verificación de software
\end{IEEEkeywords}

\section{Introducción}

En C++, un iterador, una referencia o un puntero que designa un elemento de un contenedor deja de ser utilizable cuando ciertas operaciones modifican el contenedor. La norma del lenguaje especifica, operación por operación, qué identificadores quedan invalidados \cite{iso17}, y usar uno invalidado es comportamiento indefinido (\textit{undefined behavior}, UB). La regla CTR51-CPP del estándar de codificación segura CERT del Software Engineering Institute (SEI) prohíbe expresamente ese uso \cite{cert}. El defecto es difícil de ver en una revisión: el programa suele funcionar mientras la memoria del elemento no se reutilice, y falla después, lejos de la causa.

La dificultad no está en la idea, sino en el detalle de las reglas. Un \code{push\_back} sobre \code{std::vector} invalida todo solo si se agota la capacidad; el mismo \code{push\_back} sobre \code{std::list} no invalida nada; \code{erase} sobre un \code{vector} invalida el elemento borrado y los posteriores, pero conserva los anteriores; y una inserción en el frente de un \code{std::deque} invalida los iteradores pero conserva las referencias. Decidir la validez exige saber a qué contenedor pertenece cada identificador, qué categoría tiene y dónde está su elemento respecto del punto de la mutación. Ese seguimiento es indecidible en general \cite{landi92}, \cite{ramalingam94}, de modo que toda herramienta práctica aproxima \cite{livshits15}. La forma de aproximar importa: los desarrolladores abandonan los analizadores que producen falsas alarmas \cite{christakis16}, \cite{guo23}.

Clang Static Analyzer (CSA) explora los caminos de ejecución de cada función mediante ejecución simbólica \cite{king76} y admite extensiones, llamadas \textit{checkers} (término que se conserva en inglés), que observan los eventos de cada camino \cite{csadev}. Sus checkers de iteradores, todavía experimentales, siguen iteradores pero no referencias ni punteros a elementos \cite{csadoc}, \cite{balogh18}. Otras herramientas tratan una lista de operaciones como siempre invalidantes \cite{saxena26}, \cite{infer}, \cite{cppcheck} o se limitan a \code{std::vector} \cite{infer}, \cite{liu25}. Esas herramientas tratan como siempre invalidante una inserción que solo invalida al agotarse la capacidad \cite{saxena26}, \cite{infer}, \cite{cppcheck} o la modelan solo para \code{std::vector} \cite{liu25}. Además, la evaluación más cercana usa un banco construido a mano \cite{liu25}, y los bancos generados para evaluar analizadores obtienen sus etiquetas de los defectos inyectados o de la comparación entre analizadores, no de la norma del lenguaje \cite{kashyap19}, \cite{klinger19}.

El objetivo de este trabajo es diseñar, implementar y evaluar un checker para CSA que modele las reglas de invalidación de iteradores, referencias y punteros en contenedores seleccionados de la biblioteca estándar, y comparar su precisión, exhaustividad y medida F1 con la de la configuración de referencia del analizador. Los aportes son:
\begin{itemize}
\item un modelo de invalidación que consulta las reglas de C++17 escritas como datos, con su cláusula, según el contenedor, la operación, la categoría del identificador y la posición relativa del elemento;
\item el checker \code{ucsm.ContainerInvalidation}, que se carga como \textit{plugin} en CSA, cubre diez contenedores y distingue avisos confirmados de avisos posibles, y reconoce las garantías que da el código;
\item un banco de 750 casos etiquetados a partir de la misma norma, validado con dos instrumentos dinámicos, y un protocolo que atribuye cada aviso a su sitio y separa los incidentes;
\item una evaluación comparativa en el banco y en cuatro proyectos de código abierto.
\end{itemize}

La Sección~II revisa los trabajos relacionados. La Sección~III describe el procedimiento, el desarrollo del checker y el diseño experimental. La Sección~IV presenta los resultados, la Sección~V los discute y la Sección~VI concluye.

\section{Trabajos relacionados}

\subsection{Herramientas de análisis estático}

Los checkers experimentales de CSA para iteradores, entre ellos \code{InvalidatedIterator}, \code{IteratorModeling} y \code{ContainerModeling}, asocian cada iterador a su contenedor y lo marcan como inválido ante ciertas operaciones \cite{csadoc}, \cite{balogh18}. Están desactivados por omisión y no siguen referencias ni punteros a elementos; el checker \code{cplusplus.InnerPointer} cubre punteros internos, pero solo de \code{std::string} \cite{csadoc}.

El análisis de tiempos de vida de Clang incorporó en febrero de 2026 la detección de usos tras una invalidación de contenedores de la biblioteca estándar \cite{saxena26}. Es un análisis de flujo de datos intraprocedural sobre préstamos: cualquier llamada de una lista de métodos invalida todos los préstamos del contenedor, sin considerar la capacidad ni la posición, y el caso \code{int\& r = v[0]; v.push\_back(4);} figura todavía como no detectado. La propuesta de perfil de tiempos de vida para C++ adopta también una regla conservadora sobre el dueño \cite{sutter19}, y el verificador de préstamos de Rust rechaza cualquier mutación mientras exista un préstamo vivo \cite{rust23}, aunque una reserva previa la hiciera segura.

Infer informa \code{VECTOR\_INVALIDATION} para direcciones dentro de un \code{std::vector} e incluye \code{push\_back}, \code{reserve} y \code{shrink\_to\_fit} entre las operaciones que invalidan \cite{infer}. IKOS4VecErr, basado en interpretación abstracta \cite{cousot77}, sigue el tamaño y el estado de los iteradores de \code{std::vector} y supera a CSA, Infer y Cppcheck en un banco construido a mano \cite{liu25}. Cppcheck decide con dos atributos por clase de contenedor, inserción inestable y borrado inestable \cite{cppcheck}. La consulta Itergator para CodeQL busca condiciones necesarias, pero no suficientes, de invalidación y su autor reconoce un número significativo de falsos positivos \cite{higgs20}; PVS-Studio informa mutaciones dentro de un \code{for} por rango \cite{pvs}. En C++ también se han propuesto marcos dedicados a usos tras liberación de punteros \cite{teodorescu24}.

Los instrumentos dinámicos, como el modo de depuración de libstdc++ \cite{libstdcxx} y AddressSanitizer \cite{asan}, solo observan lo que ocurrió en una ejecución: un UB que no reutiliza memoria puede pasar inadvertido.

\subsection{Comparación}

La Tabla~I resume las diferencias. Las filas de otras herramientas provienen de su documentación o de su código fuente; no se ejecutaron sobre el banco de este trabajo.

\begin{table*}[!b]
\caption{Herramientas que detectan invalidación de identificadores en contenedores de C++}
\label{tab:comparacion}
\centering
\footnotesize
\setlength{\tabcolsep}{4pt}\begin{tabular}{@{}>{\raggedright\arraybackslash}p{3.5cm} >{\raggedright\arraybackslash}p{2.8cm} >{\raggedright\arraybackslash}p{2.6cm} >{\raggedright\arraybackslash}p{3.0cm} >{\raggedright\arraybackslash}p{4.4cm}@{}}
\hline
\textbf{Herramienta} & \textbf{Técnica} & \textbf{Identificadores} & \textbf{Contenedores} & \textbf{Invalidación que depende del estado}\\
\hline
CSA, \code{InvalidatedIterator} \cite{csadoc} & ejecución simbólica & iteradores & los modelados por CSA & no se distingue\\
Clang, tiempos de vida \cite{saxena26} & flujo de datos intraprocedural & referencias e iteradores & sin \code{list} ni \code{forward\_list} & siempre invalida\\
Infer, \code{VECTOR\_INVALIDATION} \cite{infer} & análisis interprocedural & direcciones de elementos & \code{vector} & siempre invalida, incluso con \code{reserve}\\
IKOS4VecErr \cite{liu25} & interpretación abstracta & iteradores & \code{vector} & modela el tamaño\\
Cppcheck \cite{cppcheck} & flujo de valores & los tres & dos atributos por clase & siempre invalida\\
Itergator, CodeQL \cite{higgs20} & consulta de flujo de datos & iteradores & biblioteca estándar y tipos propios & no se modela\\
\textbf{Este trabajo} & ejecución simbólica (CSA) & \textbf{los tres} & \textbf{diez} & \textbf{aviso posible y garantías}\\
\hline
\end{tabular}
\end{table*}

\subsection{Necesidad que aborda el checker}

Ninguna de las herramientas revisadas combina las tres propiedades que exige la norma: aplicar la regla según contenedor, categoría y posición; seguir referencias y punteros además de iteradores; y tratar por separado la invalidación que depende del estado, sin convertirla en una falsa alarma ni en silencio. Tampoco se ha publicado una evaluación cuyo banco derive sus etiquetas de la norma y las contraste con instrumentos dinámicos \cite{kashyap19}, \cite{klinger19}.

\begin{figure*}[!t]
\centering
\begin{tikzpicture}[
  font=\footnotesize,
  caja/.style={draw, rounded corners=1.5pt, align=left, inner sep=3pt, text width=#1},
  flecha/.style={-{Latex[length=1.8mm]}, thick},
  node distance=3.5mm and 3.5mm]
\node[caja=3.95cm, fill=black!8] (norma) {%
\textbf{(a) Norma C++17 como datos}\\[1pt]
cada regla con su cláusula\\[1pt]
$N(m,k,r)\in\{$\textsc{Preserva}, \textsc{Invalida}, \textsc{Depende}$(c)\}$\\
$m$: contenedor, operación, punto\\
$k$: iterador, referencia, puntero\\
$r$: antes, mismo, después, cualquiera\\
$c$: capacidad agotada o \textit{rehash}};
\node[caja=3.1cm, right=14mm of norma.north east, anchor=north west, yshift=-6mm] (obt) {%
\textbf{(b) Obtener}\\[1pt]
\code{begin()}, \code{find()}, \code{*it}, \code{c[i]}, \code{front()}, \code{insert().first}\\[1pt]
$h=(\gamma,k,p,v)$ con $v={}$\textsc{Válido}};
\node[caja=4.1cm, right=of obt.north east, anchor=north west] (mut) {%
\textbf{(c) Mutar el contenedor $\gamma$}\\[1pt]
$R$: relaciones de $p$ con el punto $q$\\
$e=N(m,k,r)$, para todo $r\in R$\\
\textsc{Preserva}: se traslada $p$\\
\textsc{Invalida}: $v \leftarrow{}$\textsc{Inválido} y nota\\
\textsc{Depende}$(c)$: con $\kappa\geq 1$, como \textsc{Preserva}; con $\kappa=0$, $v \leftarrow v \sqcup{}$\textsc{NoResuelto}};
\node[caja=3.45cm, right=of mut.north east, anchor=north west] (uso) {%
\textbf{(d) Usar}\\[1pt]
\textsc{Válido}: sin aviso\\
\textsc{NoResuelto}: aviso posible si se pide (\code{Possibly})\\
\textsc{Inválido}: aviso confirmado con la ruta obtener, mutar, usar};
\node[caja=7.75cm, fill=black!3, anchor=north west] (leyes) at ($(obt.west |- mut.south)+(0,-3mm)$) {%
\textbf{Leyes:} \textsc{Válido} $\sqsubseteq$ \textsc{NoResuelto} $\sqsubseteq$ \textsc{Inválido}; la validez solo sube. \code{reserve(size()+n)} fija $\kappa=n$; \code{it = c.erase(it)} da un identificador nuevo y válido. Si $\gamma$ escapa a código no analizado, se olvidan los válidos.};
\node[draw, dashed, rounded corners=2pt, inner sep=4pt, fit=(obt)(mut)(uso)(leyes)] (csa) {};
\node[anchor=south west, font=\footnotesize\bfseries, inner sep=1pt] at (csa.north west) {Checker en cada camino simbólico de CSA};
\node[caja=17.7cm, fill=black!8, anchor=north west] (eval) at ($(norma.west |- csa.south)+(0,-8mm)$) {%
\textbf{(e) Evaluación:} banco de casos etiquetados con la misma norma (\code{UB}, \code{SAFE}, \code{POLICY}), con etiquetas contrastadas con el modo de depuración de libstdc++ y AddressSanitizer. CSA solo frente a CSA con el checker; precisión, exhaustividad y F1 por sitio de uso.};
\draw[flecha] (norma.east |- obt.west) -- node[above=1pt, align=center, font=\footnotesize]{tabla\\generada} (obt.west);
\draw[flecha] (obt) -- (mut);
\draw[flecha] (mut) -- (uso);
\draw[flecha] (uso.south) -- node[right]{avisos} (uso.south |- eval.north);
\draw[flecha, dashed] (norma.south) -- node[right, align=left]{etiqueta: máximo de\\la validez de los\\identificadores} (norma.south |- eval.north);
\end{tikzpicture}
\caption{Modelo de invalidación. La norma decide el efecto de cada mutación según el contenedor, la categoría del identificador y su relación con el punto (a); el checker sigue cada identificador desde su obtención (b) y a través de las mutaciones (c) hasta su uso (d). La misma norma etiqueta el banco con el que se compara CSA solo con CSA y el checker (e).}
\label{fig:modelo}
\end{figure*}

\section{Metodología}

\subsection{Procedimiento general}

El trabajo siguió cinco etapas: 1) viabilidad de la integración de un checker externo en CSA; 2) transcripción de las reglas de C++17 a datos; 3) construcción del banco, validación dinámica de sus etiquetas y medición de la configuración de referencia; 4) desarrollo incremental del checker, de iteradores a referencias y punteros, garantías y rutas de diagnóstico; y 5) evaluación final con el banco congelado. La comparación es pareada: las condiciones usan el mismo banco, el mismo compilador y las mismas opciones, y solo difieren en la carga del checker.

\subsection{Desarrollo del checker}

\subsubsection{Alcance normativo}
Las reglas provienen de las cláusulas de modificadores de contenedores del borrador N4659 \cite{n4659}, equivalente a C++17 \cite{iso17}. El alcance comprende \code{vector}, \code{deque}, \code{list}, \code{forward\_list}, \code{set}, \code{multiset}, \code{map}, \code{multimap}, \code{unordered\_set} y \code{unordered\_map}, y tres categorías de identificador: iterador, referencia y puntero. Una mutación $m$ es una terna (contenedor, operación, punto), donde el punto indica si actúa en el frente, el final, el interior, todo el contenedor o una clave. La norma se escribió como una función
\begin{equation}
N(m,k,r) \in \{\textsc{Preserva},\ \textsc{Invalida},\ \textsc{Depende}(c)\}
\label{eq:norma}
\end{equation}
donde $k$ es la categoría, $r \in \{$antes, mismo, después, cualquiera$\}$ es la relación del elemento con el punto y $c$ es la condición de la que depende el efecto: capacidad agotada o redistribución de cubetas (\textit{rehash}). Cada regla cita su cláusula. Un generador verifica sus leyes, por ejemplo que toda regla tenga cláusula, que todo efecto dependiente tenga condición y que una mutación que invalida las referencias de una relación invalide también sus iteradores, y produce la tabla en C++ que compila el checker.

\subsubsection{Modelo de invalidación}
La Fig.~\ref{fig:modelo} resume el modelo. En cada camino simbólico, un identificador es una tupla $h=(\gamma,k,p,v)$ con su contenedor $\gamma$, su categoría $k$, la posición $p$ de su elemento y su validez $v$. La posición se expresa como desplazamiento desde el primero o el último elemento, desde un ancla de posición desconocida o como desconocida. La validez pertenece a la cadena
\begin{equation}
\textsc{Válido} \sqsubseteq \textsc{NoResuelto} \sqsubseteq \textsc{Inválido}
\label{eq:cadena}
\end{equation}
y solo puede subir: ninguna operación devuelve la validez a un identificador, y volver a obtenerlo crea un identificador nuevo.

Cuando una mutación $m$ con punto $q$ actúa sobre $\gamma$, el checker calcula el conjunto $R$ de relaciones compatibles con $p$ y $q$ que la norma admite para $m$ y $k$. Si todas las relaciones de $R$ dan el mismo efecto $e$, la validez nueva es
\begin{equation}
v' =
\begin{cases}
v & e=\textsc{Preserva},\\
\textsc{Inválido} & e=\textsc{Invalida},\\
v & e=\textsc{Depende}(c),\ \kappa \geq 1,\\
v \sqcup \textsc{NoResuelto} & e=\textsc{Depende}(c),\ \kappa = 0,
\end{cases}
\label{eq:aplicar}
\end{equation}
donde $\kappa$ es el crédito de inserciones garantizadas del contenedor,
y, si las relaciones de $R$ discrepan, $v'=v \sqcup \textsc{NoResuelto}$. Cuando el identificador se preserva, su posición se traslada; por ejemplo, una inserción anterior lo desplaza un lugar. \textsc{Inválido} es absorbente.

El checker no inventa la capacidad, que CSA no representa como número. En cambio, registra las garantías que el código establece: \code{c.reserve(c.size()+n)} fija $\kappa=n$: concede $n$ inserciones de un elemento sin reubicación ni \textit{rehash}, cada inserción consume una unidad y una mutación cuyo crecimiento no se puede contar borra el crédito. La otra garantía es el identificador que devuelve la propia mutación, como en \code{it = c.erase(it)}, que se registra como nuevo y válido.

Al usar un identificador, el veredicto depende solo de su validez: \textsc{Válido} no produce aviso; \textsc{NoResuelto} produce un aviso posible únicamente si la opción \code{PolicyMode} vale \code{Possibly}, y el aviso nombra la garantía que falta; \textsc{Inválido} produce un aviso confirmado con la ruta completa: obtención, mutación que invalidó y uso. Si el contenedor se pasa a código cuyo cuerpo no se analiza, los identificadores válidos o no resueltos dejan de seguirse, porque no se sabe qué les hizo esa llamada; uno inválido se conserva, porque ninguna llamada lo revive.



\subsubsection{Implementación}
El checker es una biblioteca dinámica que \code{clang} carga con \code{-Xclang -load}, sin modificar LLVM \cite{lattner04}. Está dividido en un núcleo puro que contiene la tabla normativa y las decisiones de (\ref{eq:aplicar}) y no depende de Clang; una capa de traducción que reconoce en las llamadas de CSA las obtenciones, mutaciones y usos; y los manejadores de eventos del checker, que cargan el estado del camino, invocan al núcleo y registran el resultado con notas de ruta. El núcleo se verifica con pruebas unitarias que expresan sus leyes, con medición de cobertura de líneas, y el checker con pruebas de regresión de Clang que fijan los avisos esperados. La construcción usa LLVM 21.1.8 dentro de una imagen de contenedor reproducible.

\subsection{Evaluación experimental}

\subsubsection{Banco y etiquetas}
Cada caso es una función que obtiene un identificador, ejecuta una mutación y usa el identificador. El banco combina 57 pares contenedor y operación, tres categorías de identificador retenido y cuatro formas sintácticas (\textit{recta}; \textit{con guarda} contra \code{end()}; \textit{con llamada}, en la que la mutación ocurre en otra función; y \textit{en bucle}, donde también se usa el iterador de control), lo que da 684 casos, más 66 casos de buen uso con garantía: 54 con \code{reserve} y 12 con el identificador devuelto. La etiqueta de cada caso se deduce de (\ref{eq:norma}) y es el máximo, según (\ref{eq:cadena}), de la validez de los identificadores que el caso usa: \code{UB} si alguno queda inválido, \code{POLICY} si alguno depende del estado y \code{SAFE} si todos se preservan. El banco resultante tiene 342 casos \code{UB}, 306 \code{SAFE} y 102 \code{POLICY}.

Las etiquetas se contrastaron con dos oráculos dinámicos independientes: el modo de depuración de libstdc++ \cite{libstdcxx} y AddressSanitizer con UndefinedBehaviorSanitizer \cite{asan}. Cada caso se ejecutó en un proceso aislado; los casos \code{POLICY} se ejecutaron dos veces, una forzando la condición y otra impidiéndola. Una detección en un caso \code{SAFE} se considera una contradicción que detiene la publicación del banco. Ante una contradicción, se revisa la regla de etiquetado contra la norma y se corrige antes de congelar el banco; por ejemplo, en un bucle la etiqueta debe combinar el identificador retenido con el iterador de control. Un UB sin síntoma conserva su etiqueta, porque la norma no exige que se manifieste. El banco final se selló con su resumen criptográfico SHA-256.

\subsubsection{Entorno}
Clang y CSA 21.1.8, libstdc++ 13.3 y Ubuntu 24.04, dentro de la misma imagen para todas las condiciones, con \code{-std=c++17}.

\subsubsection{Configuraciones}
La condición de referencia activa los checkers \code{core}, \code{cplusplus} y los cinco checkers experimentales de iteradores y contenedores de CSA. Las condiciones experimentales añaden el checker con \code{PolicyMode=Silent} o \code{PolicyMode=Possibly}. Una cuarta condición carga solo el checker con \code{core} y \code{cplusplus}, para aislar su aporte.

\subsubsection{Protocolo}
El arnés invoca \code{clang} una vez por caso y condición, con un tiempo límite, y conserva la salida en el formato \textit{Static Analysis Results Interchange Format} (SARIF) \cite{sarif}. Un aviso cuenta para un caso solo si se atribuye al sitio de uso que el caso declara, y una compilación fallida, una caída o un tiempo agotado se registran como incidentes, no como verdaderos negativos. El arnés verifica las versiones y los resúmenes del checker y del banco antes de ejecutar.

\subsubsection{Métricas}
Sobre los casos \code{UB} y \code{SAFE}, un aviso confirmado es un verdadero positivo (VP) en un caso \code{UB} y un falso positivo (FP) en uno \code{SAFE}; la ausencia de aviso es un falso negativo (FN) o un verdadero negativo, respectivamente. Se calculan
\begin{equation}
P=\frac{\mathit{VP}}{\mathit{VP}+\mathit{FP}},\quad R=\frac{\mathit{VP}}{\mathit{VP}+\mathit{FN}},\quad F_1=\frac{2PR}{P+R}.
\label{eq:metricas}
\end{equation}
Los casos \code{POLICY} se informan aparte, porque en ellos el acierto es un aviso posible y no uno confirmado. También se mide la exhaustividad sobre el identificador que el caso retiene, para separar las detecciones debidas al iterador de control de un bucle.

\subsubsection{Pruebas complementarias}
El checker se ejecutó con \code{PolicyMode=Possibly} sobre todas las unidades de traducción de ninja, leveldb, yaml-cpp y RocksDB, y sobre ocho defectos sembrados y dos usos correctos sembrados en esos proyectos.

\section{Resultados}

La transcripción de la norma produjo 101 mutaciones y 348 reglas. El núcleo del checker tiene 40 pruebas unitarias, que cubren el 92.5 por ciento de sus líneas, y el checker tiene cinco archivos de pruebas de regresión con 30 avisos esperados. La validación dinámica del banco sumó 852 ejecuciones sin contradicciones, después de corregir 27 etiquetas cuya regla de etiquetado no seguía la norma.

Las cuatro condiciones completaron los 750 casos, 3000 ejecuciones en total, sin incidentes. La Tabla~II presenta las métricas de (\ref{eq:metricas}).

\begin{table}[!t]
\caption{Resultados en el banco de 750 casos}
\label{tab:resultados}
\centering
\footnotesize
\setlength{\tabcolsep}{3pt}
\begin{tabular}{@{}l r r r r r r r@{}}
\hline
\textbf{Condición} & \textbf{VP} & \textbf{FP} & \textbf{FN} & $\boldsymbol{P}$ & $\boldsymbol{R}$ & $\boldsymbol{F_1}$ & \textbf{\code{POLICY}}\\
\hline
Referencia (CSA) & 80 & 0 & 262 & 1.00 & 0.23 & 0.38 & 0/102\\
Checker, \code{Silent} & 342 & 0 & 0 & 1.00 & 1.00 & 1.00 & 0/102\\
Checker, \code{Possibly} & 342 & 0 & 0 & 1.00 & 1.00 & 1.00 & 102/102\\
Solo checker, \code{Possibly} & 342 & 0 & 0 & 1.00 & 1.00 & 1.00 & 102/102\\
\hline
\multicolumn{8}{@{}p{0.95\columnwidth}@{}}{\scriptsize La columna \code{POLICY} cuenta casos con aviso posible. Los verdaderos negativos son 306 en todas las condiciones.}
\end{tabular}
\end{table}

La configuración de referencia no produce falsos positivos, pero su exhaustividad es 0.23. La Tabla~III la desglosa sobre el identificador retenido: CSA detecta 48 de 128 iteradores y ninguno de los 200 casos con referencia o puntero; 32 de sus 80 verdaderos positivos se deben al iterador de control de los bucles. Su detección también depende de la forma: 0 de 82 en la forma recta y 16 de 82 cuando una guarda compara el iterador con \code{end()}. Con el checker, la detección es completa en todas las categorías y formas.

\begin{table}[!t]
\caption{Detección sobre el identificador retenido}
\label{tab:desglose}
\centering
\footnotesize
\begin{tabular}{@{}l r r r@{}}
\hline
\textbf{Subconjunto de casos \code{UB}} & \textbf{Casos} & \textbf{CSA} & \textbf{Checker}\\
\hline
Iterador retenido & 128 & 48 & 128\\
Referencia retenida & 100 & 0 & 100\\
Puntero retenido & 100 & 0 & 100\\
\hline
Forma recta & 82 & 0 & 82\\
Forma con guarda & 82 & 16 & 82\\
Forma con llamada & 82 & 16 & 82\\
Forma en bucle & 82 & 16 & 82\\
\hline
Total & 328 & 48 & 328\\
\hline
\end{tabular}
\end{table}

Con \code{PolicyMode=Possibly}, los 102 casos \code{POLICY} reciben un aviso posible y ninguno recibe un aviso confirmado; con \code{Silent}, ninguno recibe aviso. Ninguno de los 66 casos de buen uso con garantía produjo avisos en ninguna condición. La validación dinámica detectó UB en 174 de los 342 casos \code{UB} con el modo de depuración y en 207 con AddressSanitizer, y en ningún caso \code{SAFE}.

El costo del checker es bajo. La mediana del tiempo por caso pasó de 0.650 s en la referencia a 0.652 s con el checker, un 0.3 por ciento más, y la memoria residente máxima, de 119.1 MB a 119.9 MB, un 0.7 por ciento más.

En código real se analizaron 781 unidades de traducción: 29 de ninja, 39 de leveldb, 31 de yaml-cpp y 682 de RocksDB. Una unidad de ninja no compiló fuera de su sistema de construcción y ninguna provocó una caída de \code{clang}. No hubo avisos confirmados. Hubo un aviso posible en ninja que, tras revisarlo, es un falso positivo: el motor explora una inserción que no puede ocurrir, porque el checker no modela el contenido del contenedor. Los ocho defectos sembrados recibieron el aviso esperado, siete confirmados y uno posible, y los dos usos correctos sembrados no recibieron ninguno.

\section{Discusión}

La diferencia con la referencia se concentra en lo que CSA no modela: referencias y punteros a elementos, y la forma recta, en la que el iterador nunca se compara con un límite. El checker no sustituye la ejecución simbólica de CSA, sino que la usa para seguir cada identificador por camino. Por eso hereda sus límites: dos de los defectos sembrados solo se detectaron cuando el analizador volvió a analizar como raíz las funciones que ya había expandido en línea (\code{-analyzer-inlining-mode=all}).

El aviso posible responde a la falta de información sin decidir por el usuario. Una herramienta que invalida siempre avisaría en los 66 casos de buen uso del banco, y una que no distingue el estado callaría en los 102 casos \code{POLICY}. Esta afirmación se deduce de la documentación de esas herramientas; medirla requiere ejecutarlas sobre el banco.

\textit{Amenazas a la validez.} El checker y las etiquetas derivan de la misma norma; el resultado perfecto en el banco mide la fidelidad de la implementación a la norma escrita, no la corrección de esa transcripción. Lo mitigan los dos oráculos dinámicos independientes, la revisión de las reglas de etiquetado y la evaluación en código real con defectos sembrados. El banco es sintético y sus formas son acotadas; el código real muestra ausencia de caídas y de falsas alarmas confirmadas, no exhaustividad. Los resultados valen para C++17, diez contenedores, libstdc++ 13 y el análisis dentro de cada unidad de traducción; \code{std::string} no se incluyó.

\section{Conclusiones}

Se presentó un checker para CSA que aplica las reglas de invalidación de C++17 escritas como datos, según el contenedor, la categoría del identificador y la posición del elemento, y que separa la invalidación confirmada de la que depende del estado. En un banco de 750 casos derivado de la norma y validado dinámicamente, elevó la exhaustividad de 0.23 a 1.00 sin falsos positivos y con un sobrecosto mediano de tiempo de 0.3 por ciento, y en 781 unidades de traducción de código real no produjo avisos confirmados falsos. El trabajo futuro comprende ejecutar Clang, Infer y Cppcheck sobre el mismo banco, incorporar \code{std::string} y las operaciones de C++20, y analizar entre unidades de traducción.

\IEEEtriggeratref{23}
\begin{thebibliography}{00}
\bibitem{iso17} \textit{Programming Languages---C++}, ISO/IEC Standard 14882:2017, Int. Org. Standardization, Geneva, Switzerland, Dec. 2017.
\bibitem{cert} Software Eng. Inst. ``CTR51-CPP. Use valid references, pointers, and iterators to reference elements of a container.'' SEI CERT C++ Coding Standard. Accessed: Oct. 20, 2026. [Online]. Available: \url{https://wiki.sei.cmu.edu/confluence/spaces/cplusplus/pages/88046457/CTR51-CPP.+Use+valid+references+pointers+and+iterators+to+reference+elements+of+a+container}
\bibitem{landi92} W. Landi, ``Undecidability of static analysis,'' \textit{ACM Lett. Program. Lang. Syst.}, vol. 1, no. 4, pp. 323--337, Dec. 1992, doi: 10.1145/161494.161501.
\bibitem{ramalingam94} G. Ramalingam, ``The undecidability of aliasing,'' \textit{ACM Trans. Program. Lang. Syst.}, vol. 16, no. 5, pp. 1467--1471, Sep. 1994, doi: 10.1145/186025.186041.
\bibitem{livshits15} B. Livshits \textit{et al.}, ``In defense of soundiness: A manifesto,'' \textit{Commun. ACM}, vol. 58, no. 2, pp. 44--46, Feb. 2015, doi: 10.1145/2644805.
\bibitem{christakis16} M. Christakis and C. Bird, ``What developers want and need from program analysis: An empirical study,'' in \textit{Proc. 31st IEEE/ACM Int. Conf. Automated Softw. Eng. (ASE)}, Singapore, Sep. 2016, pp. 332--343, doi: 10.1145/2970276.2970347.
\bibitem{guo23} Z. Guo \textit{et al.}, ``Mitigating false positive static analysis warnings: Progress, challenges, and opportunities,'' \textit{IEEE Trans. Softw. Eng.}, vol. 49, no. 12, pp. 5154--5188, Dec. 2023, doi: 10.1109/TSE.2023.3329667.
\bibitem{king76} J. C. King, ``Symbolic execution and program testing,'' \textit{Commun. ACM}, vol. 19, no. 7, pp. 385--394, Jul. 1976, doi: 10.1145/360248.360252.
\bibitem{csadev} LLVM Project. ``Checker developer manual.'' Clang Static Analyzer. Accessed: Oct. 20, 2026. [Online]. Available: \url{https://clang-analyzer.llvm.org/checker_dev_manual.html}
\bibitem{csadoc} LLVM Project. ``Available checkers.'' Clang Static Analyzer Documentation. Accessed: Oct. 20, 2026. [Online]. Available: \url{https://clang.llvm.org/docs/analyzer/checkers.html}
\bibitem{balogh18} \'{A}. Balogh. (Apr. 2018). \textit{Finding iterator-related errors with Clang Static Analyzer}. Presented at EuroLLVM 2018, Bristol, U.K. [Online]. Available: \urlbalogh
\bibitem{saxena26} U. Saxena. \textit{[LifetimeSafety] Detect Use-After-Invalidation for STL Containers}. (Feb. 6, 2026). GitHub: llvm/llvm-project, Pull Request 179093, commit 21f3875. Accessed: Oct. 20, 2026. [Online]. Available: \url{https://github.com/llvm/llvm-project/pull/179093}
\bibitem{infer} Meta Platforms. ``List of all issue types.'' Infer. Accessed: Oct. 20, 2026. [Online]. Available: \url{https://fbinfer.com/docs/all-issue-types/}
\bibitem{cppcheck} Cppcheck Team. \textit{Cppcheck: lib/checkstl.cpp}. (Sep. 19, 2026). GitHub: danmar/cppcheck. (v2.22.0). Accessed: Oct. 20, 2026. [Online]. Available: \url{https://github.com/danmar/cppcheck/blob/2.22.0/lib/checkstl.cpp}
\bibitem{liu25} S. Liu, L. Chen, G. Fan, B. Yin, C. Huang, and J. Wang, ``Detecting vector container errors in C++ programs via abstract interpretation,'' in \textit{Proc. 26th Int. Conf. Formal Eng. Methods (ICFEM)}, in Lecture Notes Comput. Sci., vol. 16229, Hangzhou, China, Nov. 2025, pp. 117--133, doi: 10.1007/978-981-95-4213-0\_7.
\bibitem{kashyap19} V. Kashyap \textit{et al.}, ``Automated customized bug-benchmark generation,'' in \textit{Proc. IEEE 19th Int. Work. Conf. Source Code Anal. Manipulation (SCAM)}, Cleveland, OH, USA, Sep. 2019, pp. 103--114, doi: 10.1109/SCAM.2019.00020.
\bibitem{klinger19} C. Klinger, M. Christakis, and V. W\"{u}stholz, ``Differentially testing soundness and precision of program analyzers,'' in \textit{Proc. 28th ACM SIGSOFT Int. Symp. Softw. Testing Anal. (ISSTA)}, Beijing, China, Jul. 2019, pp. 239--250, doi: 10.1145/3293882.3330553.
\bibitem{sutter19} H. Sutter, ``Lifetime safety: Preventing common dangling,'' ISO/IEC JTC1/SC22/WG21, Geneva, Switzerland, Rep. P1179R1, Nov. 22, 2019. [Online]. Available: \url{https://www.open-std.org/jtc1/sc22/wg21/docs/papers/2019/p1179r1.pdf}
\bibitem{rust23} S. Klabnik and C. Nichols, \textit{The Rust Programming Language}, 2nd ed. San Francisco, CA, USA: No Starch Press, 2023.
\bibitem{cousot77} P. Cousot and R. Cousot, ``Abstract interpretation: A unified lattice model for static analysis of programs by construction or approximation of fixpoints,'' in \textit{Proc. 4th ACM SIGACT-SIGPLAN Symp. Princ. Program. Lang. (POPL)}, Los Angeles, CA, USA, Jan. 1977, pp. 238--252, doi: 10.1145/512950.512973.
\bibitem{higgs20} K. Higgs, ``Detecting iterator invalidation with CodeQL,'' \textit{Trail of Bits Blog}, Oct. 9, 2020. [Online]. Available: \url{https://blog.trailofbits.com/2020/10/09/detecting-iterator-invalidation-with-codeql/}
\bibitem{pvs} PVS-Studio. ``V789. Iterators for the container, used in the range-based for loop, become invalid upon a function call.'' PVS-Studio Documentation. Accessed: Oct. 20, 2026. [Online]. Available: \url{https://pvs-studio.com/en/docs/warnings/v789/}
\bibitem{teodorescu24} V.-A. Teodorescu and D. Lucanu, ``Static analysis framework for detecting use-after-free bugs in C++,'' \textit{Electron. Proc. Theor. Comput. Sci.}, vol. 410, pp. 99--115, Oct. 2024, doi: 10.4204/EPTCS.410.7.
\bibitem{libstdcxx} GNU Project. ``Chapter 18. Debug mode.'' The GNU C++ Library Manual. Accessed: Oct. 20, 2026. [Online]. Available: \url{https://gcc.gnu.org/onlinedocs/libstdc++/manual/debug_mode.html}
\bibitem{asan} LLVM Project. ``AddressSanitizer.'' Clang Documentation. Accessed: Oct. 20, 2026. [Online]. Available: \url{https://clang.llvm.org/docs/AddressSanitizer.html}
\bibitem{n4659} R. Smith, ``Working draft, standard for programming language C++,'' ISO/IEC JTC1/SC22/WG21, Geneva, Switzerland, Rep. N4659, Mar. 21, 2017. [Online]. Available: \url{https://www.open-std.org/jtc1/sc22/wg21/docs/papers/2017/n4659.pdf}
\bibitem{lattner04} C. Lattner and V. Adve, ``LLVM: A compilation framework for lifelong program analysis \& transformation,'' in \textit{Proc. Int. Symp. Code Gener. Optim. (CGO)}, San Jose, CA, USA, Mar. 2004, pp. 75--86, doi: 10.1109/CGO.2004.1281665.
\bibitem{sarif} \textit{Static Analysis Results Interchange Format (SARIF) Version 2.1.0}, OASIS Standard sarif-v2.1.0-os, Mar. 27, 2020. [Online]. Available: \url{https://docs.oasis-open.org/sarif/sarif/v2.1.0/os/sarif-v2.1.0-os.html}
\end{thebibliography}

\end{document}
